%%%%%%%% ICML 2026 EXAMPLE LATEX SUBMISSION FILE %%%%%%%%%%%%%%%%%

\documentclass{article}

% Recommended, but optional, packages for figures and better typesetting:
\usepackage{microtype}
\usepackage{graphicx}
\usepackage{subcaption}
\usepackage{booktabs} % for professional tables

% hyperref makes hyperlinks in the resulting PDF.
% If your build breaks (sometimes temporarily if a hyperlink spans a page)
% please comment out the following usepackage line and replace
% \usepackage{icml2026} with \usepackage[nohyperref]{icml2026} above.
\usepackage{hyperref}

%Package Add By Liying
\usepackage{multirow}

% Attempt to make hyperref and algorithmic work together better:
\newcommand{\theHalgorithm}{\arabic{algorithm}}

% Use the following line for the initial blind version submitted for review:
\usepackage{icml2026}

% For preprint, use
% \usepackage[preprint]{icml2026}

% If accepted, instead use the following line for the camera-ready submission:
% \usepackage[accepted]{icml2026}

\usepackage{amsmath}
\usepackage{amssymb}
\usepackage{mathtools}
\usepackage{amsthm}

\newcommand{\minisection}[1]{\vspace{0.0in} \noindent {\bf #1}}

% if you use cleveref..
\usepackage[capitalize,noabbrev]{cleveref}

%%%%%%%%%%%%%%%%%%%%%%%%%%%%%%%%
% THEOREMS
%%%%%%%%%%%%%%%%%%%%%%%%%%%%%%%%
\theoremstyle{plain}
\newtheorem{theorem}{Theorem}[section]
\newtheorem{proposition}[theorem]{Proposition}
\newtheorem{lemma}[theorem]{Lemma}
\newtheorem{corollary}[theorem]{Corollary}
\theoremstyle{definition}
\newtheorem{definition}[theorem]{Definition}
\newtheorem{assumption}[theorem]{Assumption}
\theoremstyle{remark}
\newtheorem{remark}[theorem]{Remark}

% Todonotes is useful during development; simply uncomment the next line
%    and comment out the line below the next line to turn off comments
%\usepackage[disable,textsize=tiny]{todonotes}
\usepackage[textsize=tiny]{todonotes}

%\newcommand{\SM}[1]{{\color{orange}{\bf SM:} #1}}
%\newcommand{\YX}[1]{{\color{brown}{\bf YX:} #1}}
\newcommand{\JW}[1]{{\color{magenta}{\bf JW:} #1}}

% The \icmltitle you define below is probably too long as a header.
% Therefore, a short form for the running title is supplied here:
\icmltitlerunning{Revisiting Sharpness in Low-rank Subspaces for Continual Learning}

\begin{document}

\twocolumn[
  \icmltitle{
  Revisiting Sharpness in Low-rank Subspaces for Continual Learning}
  % Adapter Subspace Sharpness Matters for LoRA-Based Continual Learning: A PAC-Bayesian Perspective}

  % It is OKAY to include author information, even for blind submissions: the
  % style file will automatically remove it for you unless you've provided
  % the [accepted] option to the icml2026 package.

  % List of affiliations: The first argument should be a (short) identifier you
  % will use later to specify author affiliations Academic affiliations
  % should list Department, University, City, Region, Country Industry
  % affiliations should list Company, City, Region, Country

  % You can specify symbols, otherwise they are numbered in order. Ideally, you
  % should not use this facility. Affiliations will be numbered in order of
  % appearance and this is the preferred way.
  \icmlsetsymbol{equal}{*}

  \begin{icmlauthorlist}
    \icmlauthor{Firstname1 Lastname1}{equal,yyy}
    \icmlauthor{Firstname2 Lastname2}{equal,yyy,comp}
    \icmlauthor{Firstname3 Lastname3}{comp}
    \icmlauthor{Firstname4 Lastname4}{sch}
    \icmlauthor{Firstname5 Lastname5}{yyy}
    \icmlauthor{Firstname6 Lastname6}{sch,yyy,comp}
    \icmlauthor{Firstname7 Lastname7}{comp}
    %\icmlauthor{}{sch}
    \icmlauthor{Firstname8 Lastname8}{sch}
    \icmlauthor{Firstname8 Lastname8}{yyy,comp}
    %\icmlauthor{}{sch}
    %\icmlauthor{}{sch}
  \end{icmlauthorlist}

  \icmlaffiliation{yyy}{Department of XXX, University of YYY, Location, Country}
  \icmlaffiliation{comp}{Company Name, Location, Country}
  \icmlaffiliation{sch}{School of ZZZ, Institute of WWW, Location, Country}

  \icmlcorrespondingauthor{Firstname1 Lastname1}{first1.last1@xxx.edu}
  \icmlcorrespondingauthor{Firstname2 Lastname2}{first2.last2@www.uk}

  % You may provide any keywords that you find helpful for describing your
  % paper; these are used to populate the "keywords" metadata in the PDF but
  % will not be shown in the document
  \icmlkeywords{Machine Learning, ICML}

  \vskip 0.3in
]

% this must go after the closing bracket ] following \twocolumn[ ...

% This command actually creates the footnote in the first column listing the
% affiliations and the copyright notice. The command takes one argument, which
% is text to display at the start of the footnote. The \icmlEqualContribution
% command is standard text for equal contribution. Remove it (just {}) if you
% do not need this facility.

% Use ONE of the following lines. DO NOT remove the command.
% If you have no special notice, KEEP empty braces:
\printAffiliationsAndNotice{}  % no special notice (required even if empty)
% Or, if applicable, use the standard equal contribution text:
% \printAffiliationsAndNotice{\icmlEqualContribution}

\begin{abstract}
 
\end{abstract}



\section{Additional Experimental Results}


\subsection{Raw Task-Wise $\operatorname{ACC}^{\mathrm{cls}}_t$ on ImageNet-R and ImageNet-C for Fig.~\ref{fig:robustness_imagenet_cp}}

To complement Fig.~\ref{fig:robustness_imagenet_cp}, we provide the underlying raw task-wise $\operatorname{ACC}^{\mathrm{cls}}_t$ trajectories for ImageNet-C and ImageNet-P in the tables below.
All entries are {means over repeated runs} (averaged across the same set of random seeds used for the main figure).



\begin{table*}[ht]
\centering
\caption{Raw task-wise $\operatorname{ACC}^{\mathrm{cls}}_t$ trajectories ($t=1,\dots,20$) on ImageNet-C under different optimizers (mean over repeated runs).}
\label{tab:taskwise_acc_imagenetr}
\setlength{\tabcolsep}{2.2pt}
\renewcommand{\arraystretch}{0.92}
\scriptsize
\begin{tabular}{@{}l l *{20}{r}@{}}
\toprule
\textbf{Method} & \textbf{Opt.} &
\textbf{t01} & \textbf{t02} & \textbf{t03} & \textbf{t04} & \textbf{t05} &
\textbf{t06} & \textbf{t07} & \textbf{t08} & \textbf{t09} & \textbf{t10} &
\textbf{t11} & \textbf{t12} & \textbf{t13} & \textbf{t14} & \textbf{t15} &
\textbf{t16} & \textbf{t17} & \textbf{t18} & \textbf{t19} & \textbf{t20} \\
\midrule

\multirow{3}{*}{\textbf{SeqLoRA}}
& SGD
& 92.67 & 87.84 & 87.67 & 83.42 & 82.13 & 81.50 & 79.62 & 78.83 & 78.11 & 76.47
& 75.85 & 75.28 & 74.69 & 74.24 & 74.27 & 73.48 & 72.69 & 71.72 & 70.49 & 69.77 \\
& \quad $+\mathrm{AS}^{(0)}_S$
& 92.67 & 87.67 & 87.33 & 84.42 & 83.20 & 82.39 & 80.33 & 79.29 & 79.11 & 78.20
& 77.73 & 77.00 & 76.69 & 77.00 & 77.33 & 77.08 & 76.45 & 75.94 & 75.33 & 73.93 \\
& \quad $+\mathrm{AS}^{(1)}_S$
& 95.67 & 93.67 & 89.78 & 88.84 & 87.67 & 86.61 & 86.67 & 86.33 & 84.85 & 84.14
& 82.61 & 82.17 & 81.87 & 81.60 & 81.42 & 80.44 & 80.10 & 79.96 & 80.21 & 80.03 \\

\addlinespace[2pt]

\multirow{3}{*}{\textbf{IncLoRA}}
& SGD
& 92.33 & 89.84 & 88.22 & 86.58 & 83.40 & 81.22 & 80.38 & 78.46 & 75.89 & 73.77
& 73.21 & 73.61 & 73.02 & 72.95 & 71.62 & 70.17 & 68.82 & 66.30 & 64.86 & 61.77 \\
& \quad $+\mathrm{AS}^{(0)}_S$
& 92.33 & 90.83 & 90.45 & 85.92 & 87.00 & 83.39 & 81.81 & 79.75 & 78.22 & 77.87
& 76.79 & 73.47 & 72.03 & 73.43 & 71.07 & 71.44 & 70.78 & 69.63 & 67.96 & 67.98 \\
& \quad $+\mathrm{AS}^{(1)}_S$
& 96.00 & 92.17 & 88.33 & 86.92 & 85.87 & 85.11 & 86.33 & 85.17 & 84.78 & 84.03
& 82.85 & 82.56 & 81.54 & 81.52 & 80.64 & 79.73 & 79.59 & 80.32 & 79.58 & 79.47 \\

\addlinespace[2pt]

\multirow{3}{*}{\textbf{OLoRA}}
& SGD
& 92.33 & 88.50 & 87.78 & 85.59 & 83.00 & 81.28 & 80.57 & 78.58 & 75.67 & 74.10
& 73.55 & 73.36 & 73.46 & 73.33 & 71.87 & 71.40 & 70.57 & 68.33 & 67.44 & 65.12 \\
& \quad $+\mathrm{AS}^{(0)}_S$
& 92.33 & 88.67 & 90.22 & 85.67 & 86.33 & 83.00 & 81.48 & 79.83 & 78.26 & 78.33
& 77.00 & 73.97 & 71.77 & 72.50 & 70.51 & 70.36 & 69.33 & 68.54 & 66.21 & 66.98 \\
& \quad $+\mathrm{AS}^{(1)}_S$
& 96.00 & 91.17 & 88.00 & 86.42 & 85.47 & 85.83 & 86.33 & 84.96 & 84.19 & 83.07
& 82.30 & 81.53 & 81.46 & 80.69 & 80.27 & 78.75 & 79.14 & 78.98 & 79.28 & 79.35 \\

\bottomrule
\end{tabular}
\end{table*}


\begin{table*}[ht]
\centering
\caption{Raw task-wise $\operatorname{ACC}^{\mathrm{cls}}_t$ trajectories ($t=1,\dots,20$) on ImageNet-P under different optimizers (mean over repeated runs).}
\label{tab:taskwise_acc_new}
\setlength{\tabcolsep}{2.2pt}
\renewcommand{\arraystretch}{0.92}
\scriptsize
\begin{tabular}{@{}l l *{20}{r}@{}}
\toprule
\textbf{Method} & \textbf{Opt.} &
\textbf{t01} & \textbf{t02} & \textbf{t03} & \textbf{t04} & \textbf{t05} &
\textbf{t06} & \textbf{t07} & \textbf{t08} & \textbf{t09} & \textbf{t10} &
\textbf{t11} & \textbf{t12} & \textbf{t13} & \textbf{t14} & \textbf{t15} &
\textbf{t16} & \textbf{t17} & \textbf{t18} & \textbf{t19} & \textbf{t20} \\
\midrule

\multirow{3}{*}{\textbf{SeqLoRA}}
& SGD
& 96.17 & 93.42 & 92.00 & 89.88 & 88.17 & 87.86 & 86.43 & 84.92 & 85.57 & 84.63
& 83.88 & 83.71 & 82.40 & 82.02 & 83.08 & 81.95 & 80.09 & 80.17 & 78.01 & 72.76 \\
& \quad $+\mathrm{AS}^{(0)}_S$
& 95.83 & 92.67 & 92.28 & 89.54 & 88.47 & 88.58 & 87.90 & 87.29 & 87.63 & 86.82
& 85.80 & 85.60 & 84.57 & 84.21 & 84.51 & 83.58 & 83.30 & 83.70 & 82.32 & 80.95 \\
& \quad $+\mathrm{AS}^{(1)}_S$
& 99.17 & 94.92 & 95.00 & 92.71 & 91.60 & 91.92 & 89.55 & 88.69 & 89.29 & 89.10
& 87.70 & 87.71 & 87.68 & 87.62 & 88.19 & 87.51 & 86.89 & 87.02 & 86.93 & 86.08 \\

\addlinespace[2pt]

\multirow{3}{*}{\textbf{IncLoRA}}
& SGD
& 95.67 & 92.50 & 91.78 & 89.71 & 88.80 & 88.08 & 85.69 & 85.29 & 84.59 & 83.25
& 80.94 & 78.58 & 77.72 & 77.79 & 74.13 & 73.24 & 71.70 & 69.29 & 68.85 & 66.92 \\
& \quad $+\mathrm{AS}^{(0)}_S$
& 95.83 & 91.58 & 92.17 & 89.62 & 89.27 & 88.28 & 86.57 & 85.31 & 85.56 & 85.33
& 83.95 & 81.81 & 80.01 & 78.26 & 78.75 & 77.08 & 77.34 & 75.79 & 76.53 & 73.75 \\
& \quad $+\mathrm{AS}^{(1)}_S$
& 98.83 & 96.41 & 95.39 & 93.33 & 91.63 & 90.97 & 89.81 & 88.83 & 88.78 & 88.63
& 87.46 & 86.27 & 86.58 & 86.15 & 86.63 & 86.42 & 86.45 & 85.54 & 85.34 & 85.09 \\

\addlinespace[2pt]

\multirow{3}{*}{\textbf{OLoRA}}
& SGD
& 95.67 & 90.00 & 90.39 & 88.33 & 88.20 & 87.61 & 85.76 & 85.02 & 84.45 & 83.37
& 79.97 & 77.76 & 76.68 & 76.64 & 72.82 & 72.54 & 70.40 & 68.16 & 67.61 & 65.55 \\
& \quad $+\mathrm{AS}^{(0)}_S$
& 95.83 & 90.17 & 91.06 & 88.46 & 88.27 & 87.86 & 86.74 & 84.96 & 85.37 & 84.95
& 83.74 & 82.61 & 81.01 & 79.58 & 78.79 & 77.21 & 76.96 & 75.30 & 75.73 & 73.32 \\
& \quad $+\mathrm{AS}^{(1)}_S$
& 98.83 & 95.33 & 94.39 & 91.46 & 90.17 & 90.14 & 88.60 & 87.48 & 87.02 & 87.23
& 86.23 & 85.47 & 85.99 & 85.79 & 86.19 & 86.23 & 86.01 & 85.29 & 84.61 & 84.47 \\
\bottomrule
\end{tabular}
\end{table*}





\subsection{Computational Cost Analysis}


Table~\ref{tab:lora_opt_metric} reports the average per-task training time on ImageNet-R.
Across LoRA variants, adding sharpness regularization increases the computational cost relative to SGD, with $AS^{(1)}_S$ being consistently the most expensive due to the additional gradient-norm maximization, while $AS^{(0)}_S$ shows a moderate overhead.
In contrast, $RS^{(0)}_S$ incurs a smaller and more variable overhead, and can be comparable to SGD in some settings.




\end{document}

